% !TEX program = lualatex
\documentclass[a4paper,11pt]{scrartcl}

\usepackage[
  lang=en,
  mode=journal,
  theme=ocean,
  fontset=lm,
  extra-math=true,
  toc=true,
  toc-layout=page
]{synaptic}
\usepackage{booktabs,tabularx}
\newcolumntype{Y}{>{\raggedright\arraybackslash}X}

\SynapticNewTheorem[remark]{observation}{Observation}
\SynapticHeader{Synaptic User Manual}
\SynapticSubHeader{Reliable academic typesetting with LuaLaTeX}
\SynapticPreprint{v2.2.0}
\SynapticTitle{The synaptic Package: User Manual}
\SynapticAuthor{synaptic project}
\SynapticAffiliation{Open-source LaTeX community}
\SynapticEmail{synaptic@example.com}
\SynapticSubject{User manual for synaptic version 2.2.0}
\SynapticKeywords{synaptic, LaTeX3, LuaLaTeX, academic typesetting}

\begin{abstract}
  \textbf{synaptic} is a modular LaTeX3 design system for journal articles,
  books, lecture material, and compact notes. This manual describes its
  requirements, safe configuration model, metadata API, theorem system, and
  mode-specific features.
\end{abstract}

\begin{document}
\SynapticMakeTitle

\section{Getting started}

The package requires LuaLaTeX and a KOMA-Script class. The repository ships
generated package files, so a local checkout can be used without installation:

\begin{verbatim}
export TEXINPUTS=/path/to/synaptic/tex/latex/synaptic//:$TEXINPUTS
lualatex document.tex
\end{verbatim}

A minimal article is:

\begin{verbatim}
\documentclass{scrartcl}
\usepackage[mode=journal]{synaptic}
\SynapticTitle{A Short Article}
\SynapticAuthor{Ada Author}
\begin{abstract}Summary.\end{abstract}
\begin{document}
\SynapticMakeTitle
\section{Introduction}
Text.
\end{document}
\end{verbatim}

Use \texttt{scrbook} for book mode. The selected class controls the paper size;
synaptic does not silently force A4.

\section{Configuration}

\noindent\begin{tabularx}{\textwidth}{@{}l l Y@{}}
  \toprule
  \textbf{Key} & \textbf{Default} & \textbf{Values and purpose} \\
  \midrule
  \texttt{mode} & journal & journal, book, lecture, or notes \\
  \texttt{theme} & ocean & ocean, graphite, forest, midnight, or paper \\
  \texttt{lang} & en & en or zh; selects labels and CJK support \\
  \texttt{fontset} & auto & auto, xcharter, libertinus, stix2, or lm \\
  \texttt{toc} & false & whether the title workflow prints a TOC \\
  \texttt{toc-layout} & top & top, inline, or page \\
  \texttt{title-layout} & auto & inline, compact, page, cover, or mode default \\
  \texttt{numbering} & auto & section, chapter, continuous, none, or mode default \\
  \texttt{box-numbering} & shared & shared, separate, or none \\
  \texttt{color-mode} & screen & screen, print, or mono \\
  \texttt{density} & balanced & compact, balanced, or airy \\
  \texttt{measure} & auto & narrow, standard, wide, or mode default \\
  \texttt{binding-offset} & 0pt & additional book binding correction \\
  \texttt{extra-math} & false & convenience Unicode math alphabets \\
  \texttt{fine-breaking} & true & moderate emergency stretch and tolerance \\
  \bottomrule
\end{tabularx}

Only \texttt{theme}, \texttt{toc}, and \texttt{toc-layout} are safe runtime
settings:

\begin{verbatim}
\SynapticSetup{theme=paper,toc=false,toc-layout=page}
\end{verbatim}

Mode, language, font, and load-time typography options are frozen after package
loading. Attempts to change them issue a warning instead of producing a
half-reconfigured document.

\subsection{Modes}

\noindent\begin{tabularx}{\textwidth}{@{}l l Y@{}}
  \toprule
  \textbf{Mode} & \textbf{Class} & \textbf{Purpose} \\
  \midrule
  journal & scrartcl & Inline article title, metadata, running heads, statements \\
  book & scrbook & Book title sequence, parts, chapters, duplex navigation \\
  lecture & scrartcl/scrreprt & Course masthead and semantic teaching system \\
  notes & scrartcl & Compact masthead and lightweight knowledge cards \\
  \bottomrule
\end{tabularx}

\subsection{Fonts}

The \texttt{auto} setting chooses the first \emph{complete} installed bundle:
XCharter, Libertinus, STIX Two, then Latin Modern. An explicit choice reports a
clear error if a required companion font is absent. Chinese documents prefer
Noto Serif/Sans CJK SC and fall back to the Fandol fonts distributed with TeX
Live.

\section{Title and PDF metadata}

All four modes share the metadata model and render it differently:

\noindent\begin{tabularx}{\textwidth}{@{}l Y@{}}
  \toprule
  \textbf{Command} & \textbf{Meaning} \\
  \midrule
  \SynapticCmd{SynapticTitle} & Required printed title and PDF title \\
  \SynapticCmd{SynapticSubtitle} & Optional display subtitle \\
  \SynapticCmd{SynapticShortTitle} & Running-head title \\
  \SynapticCmd{SynapticDate} & Explicit date \\
  \SynapticCmd{SynapticAuthor} & Repeatable author and optional affiliation mark \\
  \SynapticCmd{SynapticAffiliation} & Repeatable affiliation \\
  \SynapticCmd{SynapticEmail} & Repeatable clickable contact address \\
  \SynapticCmd{SynapticSubject} & PDF subject \\
  \SynapticCmd{SynapticKeywords} & Printed and PDF keywords \\
  \SynapticCmd{SynapticHeader} & Submission header override \\
  \SynapticCmd{SynapticSubHeader} & Secondary header \\
  \SynapticCmd{SynapticPreprint} & Preprint identifier \\
  \SynapticCmd{SynapticArxiv} & Linked arXiv identifier \\
  \SynapticCmd{SynapticDedicated} & Dedication \\
  \SynapticCmd{SynapticMakeTitle} & Validate and render all title data \\
  \bottomrule
\end{tabularx}

Journal mode accepts publication information, correspondence, and author notes;
book mode accepts half-title, publisher, edition, ISBN, and copyright fields;
lecture mode accepts course, code, lecture number, instructor, and semester.
Use \verb|\SynapticStatement{heading}{body}| for journal funding, data, ethics,
or conflict statements; the starred form omits its table-of-contents entry.
Standard \verb|\title|, \verb|\author|, and \verb|\date| are imported when the
Synaptic field is empty. PDF metadata receives one Unicode-safe conversion.

\section{Theorems and mathematics}

The built-in \texttt{theorem}, \texttt{lemma}, \texttt{proposition}, and
\texttt{definition} environments share one counter. Books use chapter scope by
default; shorter modes use sections. Definitions use upright body text;
proposition-style statements use italics in English and upright CJK text.

\begin{theorem}[Gaussian integral]
  The integral satisfies
  \[
    \int_0^\infty e^{-x^2}\,\mathrm{d}x=\frac{\sqrt{\pi}}{2}.
  \]
\end{theorem}

\begin{proof}
  Square the integral, change to polar coordinates, and evaluate the radial
  integral.
\end{proof}

New theorem types accept a style:

\begin{verbatim}
\SynapticNewTheorem[remark]{observation}{Observation}
\end{verbatim}

The available styles are \texttt{plain}, \texttt{definition}, and
\texttt{remark}.

\begin{observation}
  With \texttt{extra-math=true}, commands such as \(\symbb{R}\),
  \(\symcal{F}\), and \(\symfrak{g}\) are available.
\end{observation}

\section{Themes and lecture boxes}

Themes may be switched with \verb|\SynapticUseTheme{name}|. A custom theme is
defined and activated using xcolor expressions:

\begin{verbatim}
\SynapticDefineTheme{brand}{blue!70!black}{gray!80!black}[orange]
\end{verbatim}

Lecture mode adds the collision-resistant environments \SynapticPkg{synexample},
\SynapticPkg{synremark}, \SynapticPkg{synwarning}, \SynapticPkg{synexercise},
and \SynapticPkg{synbox}. Examples and exercises share the theorem counter by
default; remarks and warnings are unnumbered. The first optional argument is a
title and the second is a label. The generic v2.0 aliases are created only when
their names are free. Course units may also use \texttt{synlearninggoals} and
\texttt{synlecturesummary}.

\section{Book workflow and acknowledgments}

Use \SynapticCmd{SynapticFrontMatter}, \SynapticCmd{SynapticMainMatter}, and
\SynapticCmd{SynapticBackMatter} around the corresponding parts of a book.
\SynapticCmd{SynapticBookContents}, chapter subtitles, epigraphs, and
list/appendix helpers complete the long-form workflow. Notes mode provides
\texttt{synkeyidea}, \texttt{synquestion}, \texttt{syntodo}, and
\texttt{synsummary}; class-level two-column layouts are supported. The
\SynapticCmd{SynapticAcknowledgments} command creates a chapter in book mode and
a section in other modes. Its starred form suppresses the TOC entry.

\section{Development commands}

\begin{verbatim}
l3build unpack   # regenerate .sty and .def files from synaptic.dtx
l3build check    # run the LuaLaTeX regression suite
l3build doc      # build the documented source and manuals
l3build ctan     # create release archives
\end{verbatim}

\SynapticAcknowledgments*
The project acknowledges the LaTeX3, KOMA-Script, fontspec, and broader TeX
communities.

\end{document}
